\documentclass{article}

\usepackage{fancyhdr}
\usepackage{extramarks}
\usepackage{amsmath}
\usepackage{amsthm}
\usepackage{amsfonts}
\usepackage{tikz}
\usepackage[plain]{algorithm}
\usepackage{algpseudocode}
\usepackage{listings}
\usepackage{xcolor}

\definecolor{codegreen}{rgb}{0,0.6,0}
\definecolor{codegray}{rgb}{0.5,0.5,0.5}
\definecolor{codepurple}{rgb}{0.58,0,0.82}
\definecolor{backcolour}{rgb}{0.95,0.95,0.92}

\lstset{
	backgroundcolor=\color{backcolour},
	commentstyle=\color{codegreen},
	keywordstyle=\color{magenta},
	numberstyle=\tiny\color{codegray},
	stringstyle=\color{codepurple},
	basicstyle=\ttfamily\footnotesize,
	breakatwhitespace=false,
	breaklines=true,
	captionpos=b,
	keepspaces=true,
	numbers=left,
	numbersep=5pt,
	showspaces=false,
	showstringspaces=false,
	showtabs=false,
	tabsize=2
}

\usetikzlibrary{automata,positioning}

%
% Basic Document Settings
%

\topmargin=-0.45in
\evensidemargin=0in
\oddsidemargin=0in
\textwidth=6.5in
\textheight=9.0in
\headsep=0.25in

\linespread{1.1}

\pagestyle{fancy}
\lhead{\hmwkAuthorName}
\chead{\hmwkClass\ (\hmwkClassInstructor\ \hmwkClassTime): \hmwkTitle}
\rhead{\firstxmark}
\lfoot{\lastxmark}
\cfoot{\thepage}

\renewcommand\headrulewidth{0.4pt}
\renewcommand\footrulewidth{0.4pt}

\setlength\parindent{0pt}

%
% Create Problem Sections
%

\newcommand{\enterProblemHeader}[1]{
	\nobreak\extramarks{}{Problem \arabic{#1} continued on next page\ldots}\nobreak{}
	\nobreak\extramarks{Problem \arabic{#1} (continued)}{Problem \arabic{#1} continued on next page\ldots}\nobreak{}
}

\newcommand{\exitProblemHeader}[1]{
	\nobreak\extramarks{Problem \arabic{#1} (continued)}{Problem \arabic{#1} continued on next page\ldots}\nobreak{}
	\stepcounter{#1}
	\nobreak\extramarks{Problem \arabic{#1}}{}\nobreak{}
}

\setcounter{secnumdepth}{0}
\newcounter{partCounter}
\newcounter{homeworkProblemCounter}
\setcounter{homeworkProblemCounter}{1}
\nobreak\extramarks{Problem \arabic{homeworkProblemCounter}}{}\nobreak{}

%
% Homework Problem Environment
%
% This environment takes an optional argument. When given, it will adjust the
% problem counter. This is useful for when the problems given for your
% assignment aren't sequential. See the last 3 problems of this template for an
% example.
%
\newenvironment{homeworkProblem}[1][-1]{
	\ifnum#1>0
	\setcounter{homeworkProblemCounter}{#1}
	\fi
	\section{Problem \arabic{homeworkProblemCounter}}
	\setcounter{partCounter}{1}
	\enterProblemHeader{homeworkProblemCounter}
}{
	\exitProblemHeader{homeworkProblemCounter}
}

%
% Homework Details
%   - Title
%   - Due date
%   - Class
%   - Section/Time
%   - Instructor
%   - Author
%

\newcommand{\hmwkTitle}{Labwork 3}
\newcommand{\hmwkDueDate}{DUE DATE}
\newcommand{\hmwkClass}{Deep Learning}
\newcommand{\hmwkClassTime}{}
\newcommand{\hmwkClassInstructor}{SonTG}
\newcommand{\hmwkAuthorName}{\textbf{Nguyen Trong Minh}}

%
% Title Page
%

\title{
	\vspace{2in}
	\textmd{\textbf{\hmwkClass:\ \hmwkTitle}}\\
	%\normalsize\vspace{0.1in}\small{Due\ on\ \hmwkDueDate}\\
	\vspace{0.1in}\large{\textit{\hmwkClassInstructor\ \hmwkClassTime}}
	\vspace{3in}
}

\author{\hmwkAuthorName}
\date{}

\renewcommand{\part}[1]{\textbf{\large Part \Alph{partCounter}}\stepcounter{partCounter}\\}

%
% Various Helper Commands
%

% Useful for algorithms
\newcommand{\alg}[1]{\textsc{\bfseries \footnotesize #1}}

% For derivatives
\newcommand{\deriv}[1]{\frac{\mathrm{d}}{\mathrm{d}x} (#1)}

% For partial derivatives
\newcommand{\pderiv}[2]{\frac{\partial}{\partial #1} (#2)}

% Integral dx
\newcommand{\dx}{\mathrm{d}x}

% Alias for the Solution section header
\newcommand{\solution}{\textbf{\large Solution}}

% Probability commands: Expectation, Variance, Covariance, Bias
\newcommand{\E}{\mathrm{E}}
\newcommand{\Var}{\mathrm{Var}}
\newcommand{\Cov}{\mathrm{Cov}}
\newcommand{\Bias}{\mathrm{Bias}}

\begin{document}
	
	\maketitle
	\pagebreak
	
	\begin{homeworkProblem}
		\solution
		
		\subsection{Algorithm Implementation}
		We model the estimator model as : 
        $$ f_{\theta}(x) = \theta_0 * x_0 + \theta_1 * x_1 + \theta_2$$
		
		\begin{lstlisting}[language=Python, caption=Foward pass]
		def forward(x:list[float] , w:list[float]): 
			return [x_ * w[0] + w[1] for x_ in x]
		\end{lstlisting}
		
		Next is a Loss function so that the model can learn from that: 
		$$ \mathcal{L}_i = - ( y_i * \log\sigma(\hat{y}_i) + (1-y_i)*\log(1-\sigma(\hat{y}_i)))  $$
		
		You might notice that we use $\sigma()$ function in the loss instead of the forward pass.
		\begin{lstlisting}[language=Python, caption={Loss function}]
		def loss(x, y, w, z) : 
			y_pred = forward(x,y,w)
			L = sum([- (z_ * math.log(sigmoid(y_p)) + (1-z_) * math.log(1-sigmoid(y_p))) for y_p, z_ in zip(y_pred, z)])
			return L / len(x)
		\end{lstlisting}
		
		The last element we need is gradient step for each $\theta_i$:
		$$ 
		\frac{\partial \mathcal{L}}{\partial \theta_0}
		=
		\frac{1}{N}
		\sum_{i=1}^{N}
		\left(
		-z_i x_i
		+
		x_i \left(1-\sigma(-\hat{y}_i)\right)
		\right)
		 $$
		 
		$$
		\frac{\partial \mathcal{L}}{\partial \theta_1}
		=
		\frac{1}{N}
		\sum_{i=1}^{N}
		\left(
		-z_i y_i
		+
		y_i \left(1-\sigma(-\hat{y}_i)\right)
		\right)
		$$
		
		$$
		\frac{\partial \mathcal{L}}{\partial \theta_2}
		=
		\frac{1}{N}
		\sum_{i=1}^{N}
		\left(
		1 - z_i - \sigma(-\hat{y}_i)
		\right)
		$$
		
		\begin{lstlisting}[language=Python, caption={Gradient calculation}]
		def gradient(x, y, w, z): 
			y_pred = forward(x, y , w)
			grad_0 = sum([-z_ * x_ + x_ * (1 - sigmoid(-y_p) ) for x_, z_ , y_p in zip(x,z,y_pred)]) / len(x)
			
			grad_1 = sum([-z_ * y_ + y_ * (1 - sigmoid(-y_p)) for y_, z_, y_p in zip(y,z,y_pred)]) / len(x)
			
			grad_2 = sum([1 - z_ - sigmoid(-y_p) for z_, y_p in zip(z, y_pred)]) / len(x)
			
			return grad_0, grad_1, grad_2 
		\end{lstlisting}
		
		Now we can combine them all, and create a training loop that update $\theta_i = \theta_i - \frac{\partial \mathcal{L}}{\partial \theta_i} * \mu$ where $\mu$ is a learning rate. 
		
		\begin{lstlisting}
		def backward(x,y,z, init=[0,1,2],lr=10**(-1),iter=10**(3) ):
			for i in range(iter): 
			L = loss(x,y,init,z)
			g_0, g_1, g_2 = gradient(x,y,init,z)
			init[0] -= g_0 * lr
			init[1] -= g_1 * lr
			init[2] -= g_2 * lr
			
			print(i, L , init)
			
			return init
		\end{lstlisting}
		
		\subsection{Learning rate effect}
		With the full algorithm for logistic regression, we can move on to the next part and investigate the effect of the learning rate to the convergence of the model.
		
		\begin{itemize}
			\item With $\mu=1$, the model failed to converge, and fluctuate its line.
			\item With $\mu <= 0.1$, the model successfully converged. The convergence speed is directly proportional to $\mu$
		\end{itemize}
		
	\end{homeworkProblem}
	
\end{document}
